\chapter{Das persönliche Command-System}
Die folgenden Befehle bilden eine kleine, wiederverwendbare Arbeitssprache.
\begin{longtable}{@{}p{.22\textwidth}p{.7\textwidth}@{}}\toprule
\textbf{Command}&\textbf{Standardverhalten}\\\midrule\endhead
\texttt{/analyse}&Fakten, Annahmen, Ursachen, Zusammenhänge und Lösungen trennen.\\
\texttt{/entscheidung}&Kriterien, Optionen, Risiken und Folgewirkungen bewerten; Empfehlung geben.\\
\texttt{/vergleich}&Alternativen nach sinnvollen Kriterien vergleichen.\\
\texttt{/plan}&Ziel in realistische Schritte, Abhängigkeiten und nächste Aktionen zerlegen.\\
\texttt{/recherche}&Aktuelle Informationen mit belastbaren Quellen untersuchen.\\
\texttt{/kritik}&Denkfehler, Lücken, Risiken und bessere Alternativen suchen.\\
\texttt{/code}&Ausführbaren, wartbaren Code entwerfen, testen und erklären.\\
\texttt{/sysml}&Anforderungen, Struktur, Verhalten, Schnittstellen und Verifikation modellieren.\\
\texttt{/unterricht}&Lernziele, Einstieg, Erklärung, Übung und Musterlösung erstellen.\\
\texttt{/dokument}&Material logisch ordnen und als professionelles Dokument ausarbeiten.\\
\texttt{/todo}&Aufgaben, Anlass, Priorität, Frist und nächsten Schritt extrahieren.\\
\texttt{/erklaere}&Vom Einfachen zum Komplexen erklären und Begriffe einführen.\\
\texttt{/lernen}&Praxisorientierten Lernpfad mit Übungen erstellen.\\\bottomrule
\end{longtable}
\section{Modifier}
\begin{tabularx}{\textwidth}{@{}lXlX@{}}\toprule
\texttt{+kurz}&nur Wesentliches&\texttt{+tief}&mehr Zusammenhänge\\
\texttt{+quellen}&Aussagen belegen&\texttt{+schritt}&konkrete Abfolge\\
\texttt{+praxis}&Umsetzung priorisieren&\texttt{+kritisch}&Gegenargumente suchen\\
\texttt{+tabelle}&tabellarisch ausgeben&\texttt{+datei}&passendes Artefakt erzeugen\\\bottomrule
\end{tabularx}
\begin{lstlisting}
/entscheidung +kritisch +tabelle
Soll unser Team die Wissensdatenbank selbst bauen oder eine Plattform nutzen?
\end{lstlisting}
